\medskip

\begin{itemize}[label=\textbullet]
\item Sur un axe gradué en mètres, on organise une course entre une tortue et un escargot.
\item La tortue part du point d'abscisse $x = 0$. \\
Elle se déplace vers la droite à une vitesse de $2$~mètres par minute.
\item L'escargot part du point d'abscisse $x = 12$. \\
Il se déplace vers la droite à une vitesse de $50$~centimètres par minute.
\item Les deux concurrents partent en même temps.
\end{itemize}





\begin{center}
\psset{unit=0.8cm,arrowsize=2pt 3}
\begin{pspicture}(-1,-0.4)(15,0.5)
\psline[linewidth=1.25pt]{->}(-1,0)(15,0)
\multido{\n=0+1}{15}{\psline(\n,-0.1)(\n,0) \uput[d](\n,0){\n}}
\uput[u](0,0){T}\uput[u](12,0){E}
\end{pspicture}
\end{center}

\begin{tabular}{@{} ll}
D'après les données: & \textbullet~~l'abscisse de T au bout de $x$ minutes est $2x$;\\
& \textbullet~~l'abscisse de E au bout de $x$ minutes est $12+0,5x$.
\end{tabular}

La tortue rejoint l'escargot quand $2x=12+0,5x$, soit $2x-0,5x=12$ soit $1,5x=12$ ou encore $x=\dfrac{12}{1,5}$, c'est-à-dire pour $x=8$. La distance parcourue alors est de $2\times 8=16$.

La tortue rattrape l'escargot au bout de 8 minutes, donc au point d'abscisse  16.
